\subsubsection{电击类玩具（e-stim）的安全使用}

电刺激玩具（e-stim）从医用经皮电神经刺激（TENS）改装机、专用情趣主机到紫光灯（violet wand，高压高频、电流极小），种类不少，但共享同一条铁律：\textbf{要控制的是电流的路径，而不只是强度}。

\begin{itemize}
\item \textbf{电流路径避开心脏}：两个电极不得分置于躯干两侧形成"跨胸回路"（如左臂与右胸之间、两侧乳头之间）；电极远离颈部与心前区——电流通过心脏可诱发心室颤动，这与电压大小无关，与路径有关；
\item \textbf{绝对禁忌}：安装心脏起搏器或植入式除颤器（ICD）者——体外电流可干扰设备工作，后果严重；有心律失常病史者；孕期（腹部与腰骶部不放置电极）；
\item \textbf{设备选择}：医用 TENS 本身有安全参数，但改装用于性刺激已偏离设计用途；情趣专用主机的输出波形与上限更适合；\textbf{任何涉及家用市电的自制装置都等同于致命工具}——不存在"调小一点就安全"的版本；
\item \textbf{使用细节}：电极与皮肤之间用导电凝胶保证接触均匀——干燥的金属电极直接贴皮肤会造成局部灼痛；强度从最低档开始缓慢上调；皮肤破损、炎症与痣体部位不放置电极；同一场景避免"电击加束缚"叠加——被束缚者无法在感觉异常时自行脱离，监护人的责任会陡然加重。
\end{itemize}
